\section{Faithful integer evaluation and bit complexity}
\label{sec:encoding}

\subsection{Uniform polynomial bounds}

Write $V_D(t)=\sum_j c_jt^j$ for the Jones polynomial of the input knot.
The following deliberately loose bounds suffice:
\begin{equation}\label{eq:coefficient-bounds}
 \sum_j|c_j|\le4^n,
 \qquad c_j=0\quad\text{if }|j|>2n.
\end{equation}
For completeness, the normalized state sum has $2^n$ terms. Splitting each
of the $n$ crossings increases the number of connected components by at
most one, so $c(s)\le n+1$. The coefficient norm of
$\delta^{c(s)-1}$ is $2^{c(s)-1}\le2^n$. The triangle inequality gives
the first bound. Before collecting terms, the smoothing monomial,
circle powers, and writhe correction have total absolute $A$-degree at
most $n+2n+3n=6n$. The normalized exponents are divisible by four for a
knot, giving the sharper $3n/2$ bound and hence the stated $2n$ bound.

Choose
\begin{equation}\label{eq:integer-base}
 \beta=\left\lceil\frac{2n+2}{8}\right\rceil,
 \qquad B=2^\beta,
 \qquad A=B^2,
 \qquad X=B^8.
\end{equation}
Then $X\ge4\cdot4^n$. These choices depend only on the crossing count.
The backend performs no probabilistic polynomial identity test.

\begin{lemma}[Injectivity of bounded integer evaluation]
\label{lem:balanced-injective}
An integer polynomial whose coefficients have absolute value strictly
less than $X/2$ is uniquely determined by its value at $X$, among
polynomials with coefficients in the balanced digit interval
$[-X/2,X/2)$. The same is true for a Laurent polynomial with a known
finite degree range, after multiplication by a fixed monomial.
\end{lemma}

\begin{proof}
The constant coefficient is the unique representative of the residue
modulo $X$ in the balanced interval. Subtract that coefficient and divide
by $X$, then repeat. The resulting digits determine all coefficients.
The procedure works for negative integers as well. Multiplication by a
known monomial reduces the Laurent case to the ordinary polynomial case.
\end{proof}

The coefficient bounds \eqref{eq:coefficient-bounds} lie strictly inside
the balanced interval. In particular, equality with the integer evaluation
of the polynomial one proves \emph{polynomial} identity, not merely
agreement at an arbitrary evaluation point.

\subsection{Clearing negative powers before contraction}

Set $t=\zeta B$ with $\zeta^4=-1$. A local contribution has exponent
$2(1-2s)+r_v$ in $B$, between $-4$ and $4$. Multiply every crossing
tensor by $B^4$. For each edge $e$, multiply its factor by
$B^{|b_e|}$. The two choices of arrow then have nonnegative $B$-exponents
$|b_e|\pm b_e$. Consequently the complete contraction is
\begin{equation}\label{eq:scaled-bracket}
 S_0=B^{K_0}\mathcal B_D(B^2),
 \qquad K_0=4n+L_b.
\end{equation}
Every local coefficient is a bounded sum of signed monomials
$\zeta^r B^e$ with $e\ge0$. Multiplication by such a monomial is a binary
shift of an integer and a phase/sign update. General multiplication of
two large integer coefficients is unnecessary in the sequential scan.

For normalization the implementation sets
\begin{equation}\label{eq:normalization-shift}
 K=\max(K_0,6n+4),\qquad S=B^{K-K_0}S_0.
\end{equation}
The expected scaled bracket of an unknot of the input writhe $a$ is the
integer
\begin{equation}\label{eq:unknot-scalar}
 U=(-1)^{a+1}(B^8+1)B^{K+6a-4}.
\end{equation}
Its exponent is nonnegative because $a\ge-n$ and $K\ge6n+4$.
It is nonzero. Thus $S=U$ can be tested without division and is equivalent
to $V_D=1$, using the faithful evaluation above.

The empty PD is treated separately. In this project it represents one
crossing-free unknot, whose bracket is $\delta$, rather than an empty
link with bracket one. The implementation returns its full polynomial
one even when the local state and transition caps are zero.

\subsection{Recovering every coefficient}

Let $a=\operatorname{wr}(D)$. By the normalization in
\eqref{eq:bracket-state-sum},
\begin{equation}\label{eq:integer-recovery}
 P:=X^{2n}V_D(X^{-1})
   =\frac{(-1)^{a+1}S\,B^{16n+4-K-6a}}{X+1}.
\end{equation}
If the exponent of $B$ in the numerator is negative, move that power to
the denominator. The quotient is an integer: the left side is an
ordinary integer polynomial evaluated at the integer $X$.
The implementation uses exact division and checks its remainder.

Write $P=\sum_{i=0}^{4n}d_iX^i$ in balanced digits. Equation
\eqref{eq:integer-recovery} and Lemma~\ref{lem:balanced-injective} give
\[
 c_j=d_{2n-j}.
\]
The decoder checks the degree and coefficient bounds, and agreement
between the decoded identity and the direct comparison $S=U$. An identity
query does not run the decoder. A full-polynomial query with $S=U$ returns
$\{0\mapsto1\}$ immediately after validating its normalization.
These checks protect the interface against inconsistent data; the tensor
identity is what establishes the mathematical value of $S$.

\subsection{All intermediate integers have polynomially many bits}

The largest possible total $B$-exponent in an expanded shifted tensor
monomial is at most
\[
 E=8n+2L_b.
\]
There are at most $2^n$ smoothing choices and $2^{2n}$ arrow assignments.
Using the triangle inequality, every intermediate integer coordinate
has at most
\begin{equation}\label{eq:bit-bound}
 \beta(8n+2L_b)+3n+O(1)=O(n^3)
\end{equation}
bits, by Theorem~\ref{thm:turn-construction}. Partial contractions obey
the same bound because all shifted exponents are nonnegative. Final
scaling, the normalization integer, and the division in
\eqref{eq:integer-recovery} also use $O(n^3)$-bit integers.
There are only $O(n)$ output coefficients, each of $O(n)$ bits.

The loose cubic bound is intentional. The observed $L_b$ can be much
smaller than its worst-case quadratic bound, and the implementation
reports it. Improving the cochain's $\ell^1$ norm can reduce large-integer
work without changing any topological argument.

\begin{theorem}[Uncapped full Jones complexity]
\label{thm:full-jones}
For any fixed crossing order of width $w$, the new backend computes the
complete Jones polynomial using at most $2^w$ frontier keys and
$\poly(n)2^{O(w)}$ deterministic bit time. Using the existing certified
separator ordering by default gives
\[
 \poly(n)2^{O(\sqrt n)}
\]
bit time for every classical knot diagram. If a supplied family of
orders satisfies $w=O(\log^2(n+2))$, its full Jones queries run in
$n^{O(\log n)}$ bit time.
\end{theorem}

\begin{proof}
Theorem~\ref{thm:binary-tensor} bounds the number of keys. Each crossing
has bounded arity, hence a bounded number of local transitions per
incoming key. Key extraction and all integer operations cost a polynomial
in $n$, using \eqref{eq:bit-bound}. The integer reconstruction also has
polynomial bit cost, even with schoolbook division.

A collision-free table or a balanced search tree gives
$\poly(n)2^w$ operation accounting. The actual Python implementation uses
a sparse integer-key dictionary. Even the conservative worst case in
which a lookup compares every represented key squares the state factor,
which remains $2^{O(w)}$. Thus the displayed deterministic bound does
not rely on an unproved constant-time hash-table assumption. Input edge
labels are normalized; the cost of reading arbitrary integer labels is
polynomial in their supplied bit length.

The baseline separator constructor and verifier have polynomial cost and
produce an order with $w=O(\sqrt n)$. The wrapper may keep a supplied or
greedy order only after comparing its actual frontier profile with the
certified one. It therefore has the same width guarantee. Substitution
gives the first general bound and the stated polylogarithmic-width case.
\end{proof}

\subsection{Integration and exact resource semantics}

The new module is \code{fast/fastunknot/spin_jones.py}. Its public exact
query returns the identity result, optional full Laurent coefficients,
turning certificate, selected order certificate, and work counters.
The optional \code{spin-faithful} choice exposes it through the existing
Jones command and recognition filter selector. The production default
Jones backend remains unchanged because a stronger asymptotic bound
does not imply better timings on small diagrams.

State and transition caps are validated before expensive preparation.
Crossing preparation, contraction, normalization, and decoding observe
the caller's cancellation checks. If a local cap is exhausted, the
exception records charged transitions and completed crossings but no
unfinished scalar or polynomial result. A completed order certificate
may still be retained for the exact fallback. If $V_D\ne1$, the evaluator
returns a knottedness obstruction. If $V_D=1$, recognition stores that
fact and continues to an independent recognizer.

Two large signed integers can be returned by the Python query. Serialized
obstruction evidence uses hexadecimal strings for those integers and for
polynomial coefficients, avoiding process-wide changes to decimal
conversion limits. A turning certificate alone proves the validity of
the geometric weights, not the correctness of a claimed contraction
output; the latter is checked by replay or by the independent oracles
used in the accompanying audits.
